\section{Broader Impacts}\label{sec:impacts}

This is a small-scale measurement study conducted entirely on public
datasets: FineWeb-Edu \citep{penedo2024fineweb}, a DCLM-derived educational
subset \citep{li2024datacomp}, OpenR1-Math, GSM8K \citep{cobbe2021training},
and MATH \citep{hendrycks2021measuring}. It releases no new model capability:
the artifact is a 0.6B-scale reproduction of a published architecture
\citep{yang2025qwen3} trained on 1.19B tokens against the original's 36T
(Figure~\ref{fig:repro-gap}), and every recipe component studied is already
public.

The intended impact is methodological. Evidence-gated claims -- multi-seed
confidence intervals, iso-FLOP controls, versioned evaluation suites, and
pre-registered gates -- make small-model results cheaper to trust, and
publishing controlled negative results reduces wasted replication compute for
other resource-constrained groups. The SFT response-masking null
(Table~\ref{tab:sft}), the GRPO null at 0.6B (Table~\ref{tab:rlvr};
pre-registered, and reported as directional under the study's own single-seed
cap), and the context-extension stage gated off by a step-0 diagnostic are
reported in full rather than left in the file drawer; the RLVR arms carry a
content-blind random-reward control \citep{shao2025spurious} so that the null
is a measurement, not an absence of one.

The energy footprint is that of one desk-side workstation rather than a
cluster: a single machine ran every experiment, wall-clock time is recorded
per run in the shipped logs, and the largest single run -- the 1.19B-token
pretrain -- completed in 2663.1 minutes (per-stage budgets in
Table~\ref{tab:hparams}). Honest nulls at this scale are cheap to produce but
rarely published; reporting them alongside the study's two attributed wins is
the contribution this paper aims to make a norm.
